% ------------------------------------------------------------------------------------------------
\section{64-bit rings}\label{sec:bit64}

With $\ell = 64$ the fixed-point slack is no concern (at $F = 12$ to $16$ the truncations never wrap in practice), but
the 2PC configurations did not run: the preprocessing had no 64-bit linear-layer triples, the ROT-preprocessing adders
existed for $k = 8, 16, 32$ only, and every optimization of the A2B was written for 32 bits. All of it now works at
$\ell = 64$ (hpmpc \texttt{acc5728}, ConvTriple \texttt{2d0b448}; \texttt{docs/BITLENGTH64.md}).

\paragraph{HE with $t = 2^{64}$.} SEAL's plaintext moduli have at most 60 bits, so the packed evaluator does BFV with
$t = 2^{64}$ itself: encoding $\lfloor q m / 2^{64} \rceil$ exactly as $D m + \lfloor r m / 2^{64} \rceil$ for
$q = D \cdot 2^{64} + r$ (which removes BFV's $(q \bmod t)$ error term), centered 64-bit weight shares, outputs switched
to two primes, composed by CRT and truncated, and an own decryption $\lfloor x \cdot 2^{64} / q' \rceil$. A product's
noise is a sum of $(e + \rho) w$ over up to $2^{19.4}$ weight terms per output coefficient (ResNet50's widest layers),
below $2^{78}$ except with probability $2^{-40}$. The flooding ($\Delta/16 = 2^{96}$) keeps the 32-bit parameters'
margin of 19 bits above it, so $q$ needs 165 bits (three 55-bit primes, $N = 8192$; the first version took 180 bits).
An ImageNet HE product then sends 1{,}013\,MiB, $2.1\times$ the 32-bit 479\,MiB: the query has $1.5\times$ the bits
per coefficient, the response about $1.9\times$, and $N = 8192$ costs the sparse output layout its $\sqrt{N}$ factor
(\cref{tab:polyn}). Output repacking works at 64 bits as well (a 53-bit special prime, 218 bits in all, the limit for
128-bit security): 464\,MiB, below the 32-bit default, for $2.8\times$ the HE time.

\paragraph{The adders.} The circuit generator's RCA and PPA builders produce any width; for the four-way PPA, which it
builds by hand per width, a width-generic tree (groups of three slices, blocks of up to four per level) reproduces
the 64-bit circuit and produces every other width. The a-known four-way circuit needed hand fixes at 32 bits (operand
order, a-known products in dot chains that must stay pending with their own mask shares); the generator now emits
them itself, and A2bits PPA4 takes the generated a-known circuit at 64 bits as well (\cref{sec:triadround9}): 828\,MiB
and 2.8--3.1\s less preprocessing than the AB circuit it took before, for 0--0.05\s more online time.

\paragraph{The cut as a narrow adder.} At 32 bits the cut of the redundant top $F$ slices identity-substitutes the
full-width circuits by hand: public constants on the vacant slices, gate skipping with mask retargeting, and triple
bookkeeping that the bake maps must mirror. At 64 bits the MSB adder is simply narrower: the A2B prepares the full width
as before and an adder of width $64 - F$ runs on slices $F, \dots, 63$; its MSB is the sign of the value, correct by
construction. The narrow adders of all eight MSB families are generated for $F \in \{8, 10, \dots, 20, 24\}$. At
$F = 12$ the RCA needs 51 instead of 63 AND gates per value (52 instead of 64 rounds), PPA and PPA4 19--20\% fewer
triples, and the bake's Boolean addition 51 instead of 63 rounds; \flag{A2B_DELAYED_CUT} and TS1's cut design follow unchanged.

The same construction exposed a 32-bit bug: under the A2B bake, the identity-substituted PPA and PPA4 circuits (AB
and a-known alike) flip DReLU for inputs near the cut's limit $2^{31-F}$, at a rate that grows with $|v|$ (40--45\% in
the top octave; 114 of 4,096 random inputs at $F = 5$, all above $2^{16}$), while RCA and the builds without the bake
are correct. With \flag{CUT_NARROW_32} the 32-bit PPA and PPA4 under the bake now take narrow adders too, which fixes it.
On ImageNet (A2bits, interleaved with \flag{CUT_NARROW_32=0}) the narrow PPA needs 33 fewer rounds and 15\,MiB less
online traffic (online 0.39 against 0.41\s in UC2) for 0.1--0.25\s more preprocessing. PPA4 first traded the a-known
circuit, which had no generated form, for the AB one, whose Beaver 3- and 4-tuples took 323\,MiB more OT material;
with the generated a-known narrow circuit it keeps the a-known one (\cref{sec:triadround9}). The a-known RCA, whose
generated circuits fold the a-known LSB carry, takes the narrow adder as well: 49 fewer rounds per inference.

\paragraph{The reshare bake and the arithmetic triples.} The reshare bake's maps (which slice consumes which random
multiplication and which 3-tuple) are read off the generated 64-bit and narrow circuits and checked by the generator
script; the generator's skipped input re-maskings become \texttt{zero\_add\_local}, which falls back to the real
re-masking for inputs whose mask is not baked, as in the 32-bit files. At two narrow widths the generator skips
re-maskings with 4-tuple fields that the bake cannot serve; there they stay. P1's ReLU preprocessing traffic drops by
15\% for PPA on CIFAR-10. One more 64-bit defect was silent: the elementwise arithmetic triples behind every
secret-by-secret product (max pooling) came from gemini's HE in SEAL's plaintext space, and its bit-length assertion
is compiled out in release builds, so with 64 bits they were random. Gilboa multiplication over the silent COTs fixed
them first (for bit $j$ of one factor, one COT of width $64 - j$); now gemini's CRT-batched HE product itself takes 64-bit
shares: the product of two shares plus a mask that hides it statistically ($2^{-40}$) needs $2 \cdot 64 + 1 + 40 = 169$
bits of CRT plaintext, five 34-bit primes, which keeps the plaintext scaling after the switch to the 49-bit prime. Its mask
was also too wide (the top limb kept the complement bits: 87-bit masks at 32 bits, harmless but wasteful). For 409{,}600
AB2 triples the HE product takes 0.23\s and 57\,MB against Gilboa's 1.42\s and 129\,MB (AB: 0.30\s, 114\,MB against
2.05\s, 248\,MB); \flag{ARITH_OT=1} keeps Gilboa.

\paragraph{Results.} \Cref{tab:bit64} compares the 18 all\_opt builds with \flag{COMPRESS=0} at 64 bits ($F = 12$) with
the same builds at 32 bits ($F = 5$), on flare / polynize at the round-9 code. Preprocessing traffic grows by
$1.7$--$2.1\times$ (the HE products $2.1\times$, the OT material $1.7$--$2.1\times$: the Boolean triples and tuples scale
with the slices, the COTs' messages with $\ell$) and online traffic by $2.0\times$; preprocessing takes $1.5$--$1.8\times$ as
long (PPA4 the most, whose tuples dominate), the online phase $1.5$--$1.9\times$. RCA pays the most rounds, 50 instead of 25
AND levels per ReLU (2{,}890 against 1{,}665 rounds in UC2, with the folded a-known RCA); the prefix adders add one level
(49 more rounds). Repacking takes
the HE products of UC2 A2bits RCA from 1{,}050 to 524\,MiB (1{,}320 instead of 1{,}870\,MiB in all, $1.37\times$ the 32-bit
traffic; hpmpc \texttt{4c2e446}) for 1.8\s more preprocessing. On CIFAR-10 (AdamW ResNet50, 256 images, UC2 A2bits RCA) the 64-bit builds classify
like the plaintext model: TS\{L\} 188 at $F = 12$ and 189 at $F = 16$, TE0 and TE1 the same (32 bits: TS\{L\} 159 / 169 at
$F = 5$ / 8; plaintext 189).

\begin{table}[t]
\centering\small
\begin{tabular}{@{}lrrrrr@{}}
\toprule
build & pre MiB & pre s & online MiB & online s & rounds\\
\midrule
UC1 A2bits RCA & 1{,}406 / 2{,}807 & 3.26 / 4.98 & 208 / 416 & 0.48 / 0.85 & 1769 / 2994\\
UC1 A2bits PPA & 1{,}508 / 3{,}021 & 4.25 / 6.66 & 307 / 620 & 0.45 / 0.79 & 949 / 998\\
UC1 A2bits PPA4 & 1{,}772 / 3{,}674 & 4.51 / 8.30 & 202 / 403 & 0.43 / 0.70 & 730 / 779\\
UC2 A2bits RCA & 966 / 1{,}873 & 3.26 / 4.93 & 166 / 331 & 0.40 / 0.74 & 1665 / 2890\\
UC2 A2bits PPA & 1{,}068 / 2{,}087 & 4.24 / 6.74 & 264 / 535 & 0.38 / 0.68 & 845 / 894\\
UC2 A2bits PPA4 & 1{,}333 / 2{,}740 & 4.42 / 7.98 & 159 / 318 & 0.36 / 0.62 & 626 / 675\\
UC3 A2bits RCA & 441 / 768 & 2.94 / 4.46 & 123 / 246 & 0.46 / 0.74 & 1665 / 2890\\
UC3 A2bits PPA & 543 / 983 & 3.94 / 6.04 & 222 / 450 & 0.41 / 0.73 & 845 / 894\\
UC3 A2bits PPA4 & 808 / 1{,}635 & 4.37 / 7.66 & 117 / 233 & 0.43 / 0.70 & 626 / 675\\
\addlinespace
UC1 reshared RCA & 1{,}229 / 2{,}521 & 2.14 / 3.51 & 239 / 474 & 0.54 / 0.95 & 1867 / 3092\\
UC1 reshared PPA & 1{,}389 / 2{,}744 & 3.97 / 6.06 & 349 / 676 & 0.57 / 0.96 & 1004 / 1047\\
UC1 reshared PPA4 & 2{,}150 / 4{,}226 & 5.54 / 9.24 & 233 / 459 & 0.52 / 0.91 & 779 / 828\\
UC2 reshared RCA & 790 / 1{,}587 & 2.10 / 3.43 & 197 / 389 & 0.47 / 0.82 & 1763 / 2988\\
UC2 reshared PPA & 949 / 1{,}810 & 3.97 / 6.06 & 306 / 591 & 0.49 / 0.82 & 900 / 943\\
UC2 reshared PPA4 & 1{,}710 / 3{,}292 & 5.45 / 9.20 & 190 / 374 & 0.46 / 0.71 & 675 / 724\\
UC3 reshared RCA & 267 / 487 & 1.68 / 2.80 & 154 / 304 & 0.51 / 0.91 & 1763 / 2988\\
UC3 reshared PPA & 450 / 758 & 3.21 / 5.01 & 264 / 506 & 0.50 / 0.83 & 900 / 943\\
UC3 reshared PPA4 & 1{,}102 / 2{,}068 & 5.00 / 8.31 & 148 / 289 & 0.49 / 0.79 & 675 / 724\\
\bottomrule
\end{tabular}
\caption{The all\_opt builds at 32 / 64 bits ($F = 5$ / 12; ImageNet, dummy weights, flare / polynize, hpmpc \texttt{4b43984},
the artifact configurations, median of 3 runs; P0's traffic sent and received). \flag{COMPRESS=0}; the cut as narrow adders
where the A2B bake runs (A2bits) and at 64 bits, the generated a-known PPA4 at both widths.}
\label{tab:bit64}
\end{table}

